\documentclass[11pt]{article}

\usepackage[margin=1in]{geometry}
\usepackage{amsmath,amssymb,amsthm}
\usepackage{graphicx}
\usepackage[hidelinks]{hyperref}

\newtheorem{theorem}{Theorem}
\newtheorem{lemma}{Lemma}

\newcommand{\R}{\mathbb{R}}
\newcommand{\N}{\mathbb{N}}
\newcommand{\dd}{\,d}

\title{A Coordinate Anisotropic Lift of Schindl's No-Convex-Equivalent Weight}
\author{Run fa\_banach\_001, agent\_super\_01}
\date{July 1, 2026}

\begin{document}
\maketitle

\section*{Status}

This is a partial result for the final question/goal in Schindl's paper
\cite{Schindl2510}.  The result gives an explicit anisotropic lift under the
natural coordinate-log-convex analogue of the one-dimensional condition
$(\omega_4)$.  It does not claim to settle every possible convention for
anisotropic Braun-Meise-Taylor weights.

\section*{Source Question}

The source paper constructs in Section 7 a one-dimensional weight whose
equivalence class contains no representative satisfying $(\omega_4)$, i.e. no
equivalent $\sigma$ for which $t\mapsto \sigma(e^t)$ is convex.  It then asks
for a higher-dimensional anisotropic analogue:

\begin{center}
\includegraphics[width=0.96\linewidth]{figures/open_problem_crop.png}
\end{center}

The final paragraph asks to ``transfer the construction of
$\varphi_\omega$ (resp. of $\omega$) in Section 7 to the higher dimensional case
(anisotropic setting)'' and notes that the growth-index argument is unclear in
that setting.

\section*{Definitions and Source Input}

For functions $F,G:\R^d\to[0,\infty)$ write $F\sim G$ if there is $C\ge 1$
such that
\[
  F(x)\le C G(x)+C,\qquad G(x)\le C F(x)+C\qquad (x\in\R^d).
\]
For one-variable functions this is the equivalence relation used in
\cite{Schindl2510}.

For an anisotropic function $\Sigma:\R^d\to[0,\infty)$ define the
coordinate-log-convex analogue of $(\omega_4)$ by requiring
\[
  \Phi_\Sigma(u_1,\ldots,u_d)
  :=\Sigma(e^{u_1},\ldots,e^{u_d})
\]
to be convex on $\R^d$.  This is the multivariable condition addressed here.

We use the output of Section 7 of \cite{Schindl2510}: there is a
one-dimensional function $\omega$ with the following properties.
\begin{itemize}
\item $\omega$ is equivalent to a continuous, nondecreasing, subadditive
  function $\rho:[0,\infty)\to[0,\infty)$ with $\rho(0)=0$.
\item $\rho$ satisfies $(\omega_3)$: for every $a\ge 1$ there is $b\in\R$ such
  that $\rho(t)\ge a\log(1+t)+b$ for all $t\ge 0$.
\item No one-dimensional function $\sigma\sim\rho$ has $t\mapsto\sigma(e^t)$
  convex.
\item The source construction gives uncountably many pairwise non-equivalent
  choices of such $\rho$.
\end{itemize}
The first and third bullets follow because Section 7 proves the obstruction for
$\omega$ and also proves that $\omega$ is equivalent to the continuous
subadditive representative $\kappa_\omega$; equivalence preserves the
obstruction.

\section*{Idea of the Proof}

The lift is intentionally simple: put independent copies of the bad
one-dimensional representative on the coordinate axes, but with unequal
coefficients so the resulting weight is genuinely non-radial.  Subadditivity
and growth survive coordinatewise.  If an equivalent anisotropic weight had the
coordinate-log convexity property, freezing all coordinates except the first at
$1$ would produce a one-dimensional convex representative equivalent to
$\rho$, exactly contradicting the Section 7 obstruction.  Thus no new
growth-index machinery is needed.

\begin{theorem}[coordinate anisotropic lift]
Let $d\ge 2$ and let $\rho$ be as above.  Define
\[
  \Omega(x_1,\ldots,x_d)
  := \rho(|x_1|)+2\rho(|x_2|)
     +\sum_{j=3}^d \rho(|x_j|).
\]
Then $\Omega:\R^d\to[0,\infty)$ is continuous, subadditive, and non-radial.  It
satisfies the anisotropic lower growth condition
\[
  \forall a\ge 1\;\exists b\in\R\;\forall x\in\R^d:\qquad
  \Omega(x)\ge a\log(1+|x|)+b.
\]
Also, for every $\lambda>0$,
\[
  \int_{\R^d} e^{-\lambda\Omega(x)}\dd x<\infty.
\]
Finally, there is no $\Sigma:\R^d\to[0,\infty)$ such that $\Sigma\sim\Omega$
and $\Phi_\Sigma(u)=\Sigma(e^{u_1},\ldots,e^{u_d})$ is convex on $\R^d$.
\end{theorem}

\begin{proof}
Continuity and $\Omega(0)=0$ are immediate from the corresponding properties of
$\rho$.  Since $\rho$ is subadditive and nondecreasing,
\[
  \rho(|x_j+y_j|)
  \le \rho(|x_j|+|y_j|)
  \le \rho(|x_j|)+\rho(|y_j|),
\]
and summing over the coordinates with the displayed positive coefficients gives
$\Omega(x+y)\le\Omega(x)+\Omega(y)$.

The function is not radial.  Indeed, for any $t$ with $\rho(t)>0$,
\[
  \Omega(t,0,\ldots,0)=\rho(t),
  \qquad
  \Omega(0,t,0,\ldots,0)=2\rho(t),
\]
although the two points have the same Euclidean norm.

For the lower growth, let $c_{\min}=1$.  Given $x\in\R^d$, choose a coordinate
$j$ with $|x_j|\ge |x|/\sqrt d$.  Since every coefficient in $\Omega$ is at
least $1$,
\[
  \Omega(x)\ge \rho(|x_j|)\ge \rho(|x|/\sqrt d).
\]
Given $a\ge 1$, use $(\omega_3)$ for $\rho$ to choose $b_0$ such that
$\rho(t)\ge a\log(1+t)+b_0$ for all $t\ge0$.  Since
$1+r/\sqrt d\ge (1+r)/\sqrt d$ for $r\ge0$,
\[
  \Omega(x)\ge a\log(1+|x|)-\frac a2\log d+b_0.
\]
This is the desired anisotropic $(\omega_3)$ condition.

Next fix $\lambda>0$.  The integral factors:
\[
  \int_{\R^d} e^{-\lambda\Omega(x)}\dd x
  =
  \prod_{j=1}^d \int_\R e^{-\lambda c_j\rho(|t|)}\dd t,
\]
where $(c_1,c_2,c_3,\ldots,c_d)=(1,2,1,\ldots,1)$.  For each $j$, choose
$q_j>1/(\lambda c_j)$.  By $(\omega_3)$, $\rho(t)\ge q_j\log(1+t)+b_j$.
Hence
\[
  \int_\R e^{-\lambda c_j\rho(|t|)}\dd t
  \le
  2e^{-\lambda c_j b_j}
  \int_0^\infty (1+t)^{-\lambda c_j q_j}\dd t
  <\infty.
\]
Thus the product is finite.

It remains to prove the obstruction to coordinate-log convex representatives.
Assume, toward a contradiction, that $\Sigma\sim\Omega$ and
$\Phi_\Sigma$ is convex.  Define a one-dimensional function
\[
  \sigma(t):=\Sigma(t,1,\ldots,1),\qquad t\ge0.
\]
Then
\[
  \Omega(t,1,\ldots,1)
  =
  \rho(t)+K,
\]
where $K=2\rho(1)+\sum_{j=3}^d \rho(1)$ is constant.  The equivalence
$\Sigma\sim\Omega$ therefore implies $\sigma\sim\rho$ as functions on
$[0,\infty)$.  But
\[
  s\mapsto \sigma(e^s)
  =
  \Sigma(e^s,1,\ldots,1)
  =
  \Phi_\Sigma(s,0,\ldots,0)
\]
is the restriction of a convex function to a line, hence is convex.  This
contradicts the one-dimensional Section 7 obstruction for $\rho$.
\end{proof}

\begin{theorem}[uncountably many lifts]
The construction above gives uncountably many pairwise non-equivalent
anisotropic weights $\Omega$ with the same properties.
\end{theorem}

\begin{proof}
Take the uncountable family of pairwise non-equivalent one-dimensional
representatives $\rho_\alpha$ supplied by Section 7 of \cite{Schindl2510}, and
form $\Omega_\alpha$ as above.  If $\Omega_\alpha\sim\Omega_\beta$, then
restricting to $(t,1,\ldots,1)$ gives
\[
  \rho_\alpha(t)+K_\alpha \sim \rho_\beta(t)+K_\beta,
\]
and hence $\rho_\alpha\sim\rho_\beta$, contradicting the choice of the family.
\end{proof}

\section*{Verification and Limitations}

The proof is deterministic and uses no computation beyond rendering the source
PDF and crop.  The local run indexes and cheap source indexes were searched for
\texttt{2510.05469}, \texttt{anisotropic convex},
\texttt{coordinate log convex}, \texttt{log-convex anisotropic}, and
\texttt{Beurling Bjorck anisotropic}; no existing packet for this coordinate
lift was found.

The result is intentionally scoped.  It answers the final question under the
coordinate-log-convex formulation stated above.  If the intended anisotropic
replacement for $(\omega_4)$ in another paper is weaker, stronger, or phrased
through a different transform, that convention should be checked separately.

\begin{thebibliography}{9}

\bibitem{Schindl2510}
G. Schindl.
\newblock \emph{On inclusion relations of weighted $L^p$-type spaces defined
in terms of weight function matrices}.
\newblock arXiv:2510.05469.

\end{thebibliography}

\end{document}
